
\documentclass{article}
\usepackage{amsmath}
\usepackage{graphicx}
\usepackage{hyperref}

\title{Introducción y Regulaciones del Powerlifting}
\author{Andrés López}

\begin{document}
Andrés López González

\section{Información sobre el Powerlifting}

El powerlifting, también conocido como levantamiento de potencia, es un deporte de fuerza que consta de tres levantamientos principales: sentadilla, press de banca y peso muerto. El objetivo es levantar el mayor peso posible en cada uno de estos ejercicios, sumando el total de los tres para determinar el rendimiento del atleta.

\subsection{Categorías de Peso y Edades}
Existen diversas categorías de peso corporal y rangos de edad para garantizar una competencia justa. Las categorías varían según la federación, pero generalmente incluyen divisiones masculinas y femeninas desde los 52 kg hasta más de 120 kg. En mujeres la tendencia es la siguiente clasificación: 40 kg, 44 kg, 48 kg, 52 kg, 56 kg, 60 kg, 67.5 kg, 69 kg, 76 kg, open (sin límite); en hombres por otro lado se divide como: 50 kg, 55 kg, 60 kg, 65 kg, 120 kg, open. Las categorías de edad son Sub-junior: 14-18 años;
Junior: 19-23 años; Open 23-40 años (pero cualquiera puede entrar)
Master 1: 40-49 años;
Master 2: 50-59 años;
Master 3: 60-69 años;
Master 4: 70+.

\subsection{Equipamiento: Raw y Equipado}
El powerlifting puede practicarse en dos modalidades principales:
\begin{itemize}
    \item \textbf{Raw}: Se permite el uso de rodilleras o vendas, cinturón y muñequeras, pero sin equipamiento de asistencia adicional.
    \item \textbf{Equipado}: Se permite el uso de trajes especiales de sentadilla y banca, que brindan soporte adicional y permiten levantar más peso.
\end{itemize}

Es importante señalar que cada levantamiento tiene requisitos técnicos específicos:
\begin{itemize}
    \item \textbf{Sentadilla}: Debe ejecutarse descendiendo hasta que la cadera esté por debajo de la rodilla antes de subir.
    \item \textbf{Press de banca}: El levantador debe bajar la barra hasta el pecho, hacer una pausa y luego presionar hasta la extensión total de los codos.
    \item \textbf{Peso muerto}: La barra debe levantarse desde el suelo hasta la extensión completa de la cadera y las rodillas, manteniendo el control.
\end{itemize}

\section{Comparación con Otras Disciplinas de Fuerza}

El powerlifting se diferencia de otras disciplinas como:
\begin{itemize}
    \item \textbf{Fisicoculturismo}: Se enfoca en la estética y el desarrollo muscular, no en la fuerza máxima.
    \item \textbf{Halterofilia}: Se centra en el arranque y el envión, movimientos más técnicos y explosivos.
    \item \textbf{CrossFit}: Incluye ejercicios variados con énfasis en resistencia, acondicionamiento y técnica.
\end{itemize}
Como se ha mencionado, el powerlifting se centra en la máxima potencia de los 3 levantamientos básicos

\section{Powerlifting Como Disciplina Competitiva}

El powerlifting es una disciplina estructurada con reglas claras y un enfoque en la progresión de la fuerza. Practicarlo sin un enfoque adecuado puede aumentar el riesgo de lesiones y limitar el desarrollo del atleta, así como no permitirá un correcto desarrollo en la fuerza del atleta por la falta de motivos. Participar en competencias fomenta la disciplina, el compromiso y la mejora continua así como la exposición a un jurado que dictaminará si verdaderamente estás haciendo las cosas como se debe.

\end{document}
